\section{Documento de diseño del proyecto}

El servidor usa tres piezas centrales del toolkit: \texttt{selector},
\texttt{stm} y \texttt{buffer}. El selector concentra los eventos de sockets; la
máquina de estados separa cada fase del protocolo; los buffers resuelven
lecturas y escrituras parciales.

\subsection{Estados SOCKS5}

\begin{center}
\texttt{HELLO -> AUTH -> REQUEST -> RESOLV? -> CONNECTING -> REQUEST\_WRITE -> COPY -> DONE}
\end{center}

\texttt{COPY} mantiene dos direcciones de relay: cliente a origen y origen a
cliente. Cada dirección tiene su estado de lectura/escritura; al recibir EOF en
un sentido se drena lo pendiente y se propaga \texttt{shutdown(SHUT\_WR)} al par.
Cuando una lectura produce datos, el relay realiza inmediatamente un único
intento de escritura no bloqueante sobre el socket opuesto. En el caso favorable
esto evita una vuelta adicional \texttt{pselect} $\rightarrow$ \texttt{write}; ante una escritura
parcial o \texttt{EAGAIN}, conserva los bytes pendientes en el buffer y delega el
drenaje posterior al selector. El límite de un intento por activación preserva la
equidad del event loop entre conexiones concurrentes.

\subsection{DNS y lifetime}

La resolución FQDN corre en un hilo detached que sólo ejecuta \texttt{getaddrinfo}
y notifica al selector. La notificación lleva identidad del objeto de conexión y
cleanup asociado. El proceso de apagado espera a que no haya resoluciones DNS en
vuelo antes de destruir selector y pool.

\subsection{PMC}

El PMC tiene su propio listener y su propia STM, pero comparte el mismo selector
del servidor. El cliente CLI usa I/O bloqueante por simplicidad, permitido por la
consigna, y encapsula el framing del protocolo.
